% Generated from article.md by scripts/render_article.py. Edit article.md.
\documentclass{misisarticle}
\grnti{28.23.37, 16.31.21}
\asjc{1702, 1706}
\oecd{1.02}
\title{Оценка параметрически эффективной адаптации языковой модели к татарскому языку}
\misisauthor{Горин Александр Олегович}{НИТУ «МИСиС»}{Москва, Россия}{alexgorin37@gmail.com}
\abstract{Для переноса освоенной моделью задачи на другой язык важно выяснить, где следует разместить ограниченное число обучаемых параметров и как изменение качества целевого языка связано с сохранением исходной задачи. В работе предлагается протокол сравнения адаптации одной небольшой открытой авторегрессионной модели к татарскому языку. Общий checkpoint обучается тематической классификации только на русских метках SIB-200. Его независимые копии затем адаптируются предсказанию следующего токена на отдельном неразмеченном татарском корпусе: с LoRA в нижних, верхних или всех блоках; вариант без языковой адаптации служит основным контролем, обновление эмбеддингов — дополнительным. Основное сравнение расположения LoRA проводится между нижними и верхними блоками при одинаковом ранге, числе блоков и наборе матриц. Качество переноса планируется измерять по macro-F1 и accuracy на татарском тесте, сохранение задачи — на русском тесте; отдельно учитываются языковая потеря, число параметров, объём данных и вычислительные затраты. Протокол исключает пересечения с оценочными предложениями и не использует татарские метки для настройки. Эксперимент ещё не проведён; численных выводов об эффективности методов работа пока не содержит.}
\keywords{языковые модели, татарский язык, межъязыковой перенос, LoRA, тематическая классификация, параметрически эффективная адаптация}
\begin{document}
\maketitle
\noindent\emph{Научный руководитель: И. И. Пионтковская, НИТУ «МИСиС».}\par
\section{Введение}

Предобученная языковая модель может освоить задачу по размеченным примерам одного языка, но перенос на другой язык зависит, среди прочего, от языкового охвата предобучения и токенизации \cite{pfeiffer2020mad,csaki2023efficiently,adelani2023sib}. Здесь рассматривается сценарий, в котором доступны русские метки задачи и неразмеченный татарский текст, но татарские метки нельзя использовать для настройки. Практический вопрос состоит в том, изменится ли перенос при ограниченном обновлении уже существующей модели и сохранится ли качество исходной задачи.

Исследование промежуточных представлений Llama-2, проведённое Wendler и соавт. \cite{wendler2024llamas}, мотивирует сравнение вмешательств в разных частях сети. Его наблюдения получены без языковой адаптации и не позволяют заранее предсказать результат такой проверки на татарской классификации.

Цель работы — задать воспроизводимое сравнение способов языковой адаптации на одной модели, одном целевом языке и одной задаче. Основной вопрос: различается ли перенос тематической классификации с русского на татарский, если LoRA с одинаковым рангом и набором матриц размещена в равном числе нижних либо верхних блоков? Второй вопрос: как изменение татарского качества соотносится с сохранением русского? Вариант LoRA во всех блоках показывает, полезно ли распределить сходное число обучаемых параметров по всей глубине модели, но его сравнение с локальными вариантами смешивает расположение и ранг адаптеров. Эта постановка не ищет глобальный минимум параметров. В статье рассматривается публичная предобученная модель, поэтому татарский нельзя называть языком, заведомо отсутствовавшим в её предобучении. Проверка действительно \emph{невиданного} языка требует модели с контролируемым составом данных предобучения и относится к дальнейшей программе диплома.

\section{Связанные работы}

\subsection{Язык и представления внутри модели}

Анализ Wendler и соавт. \cite{wendler2024llamas} относится к диагностике промежуточных представлений Llama-2, а не к обучению нового языка. Он мотивирует проверку расположения адаптации, но не даёт направленного прогноза о том, какие слои лучше изменить для тематической классификации татарских предложений. Авторы также обсуждают влияние токенизации на наблюдаемые ими языковые траектории \cite{wendler2024llamas}. Поэтому сравнение слоёв требует отдельного аудита токенизатора выбранной модели.

\subsection{Малое обновление параметров}

Aghajanyan и соавт. показали, что для ряда задач дообучения предобученных моделей хорошие решения достижимы в пространствах существенно меньшей размерности, чем полное пространство весов \cite{aghajanyan2020intrinsic}. Их эксперименты посвящены главным образом настройке на задачу, поэтому низкая внутренняя размерность языковой адаптации из них не следует. LoRA представляет обновление матрицы весов произведением двух обучаемых матриц малого ранга при замороженных исходных весах \cite{hu2021lora}. Это позволяет независимо выбирать ранг и место установки адаптеров. ROSA расширяет идею адаптации в подпространствах, последовательно изменяя выбранные подпространства \cite{hameed2024rosa}. В первом эксперименте ROSA не включена: основное сравнение ограничено расположением LoRA при сопоставимом параметрическом бюджете. Она остаётся кандидатом для следующего этапа после получения воспроизводимого базового результата.

\subsection{Перенос задачи и адаптация языка}

В MAD-X языковые адаптеры обучаются на неразмеченном тексте, а адаптер задачи — на размеченных данных исходного языка; при переносе языковой модуль заменяется \cite{pfeiffer2020mad}. Эта работа демонстрирует полезность разделения двух источников обучения, но использует прежде всего многоязычные энкодеры и адаптеры с другой архитектурой. Csaki и соавт. адаптировали модель к венгерскому и тайскому языкам, отдельно рассматривая эффективность токенизатора и сохранение качества на исходном языке \cite{csaki2023efficiently}. Их схема включает изменения словаря, продолжение предобучения и обучение на инструкциях, поэтому её результаты нельзя переносить как численное ожидание для LoRA на татарском. Для настоящей работы существенны два методических вывода: токенизацию нужно измерить до запуска, а качество исходного языка — контролировать после адаптации.

\subsection{Контекст небольших моделей}

Africa и соавт. изучали межъязыковой перенос распознавания именованных сущностей на небольших декодерных моделях, обученных с контролируемым отсутствием целевых языков в предобучении \cite{africa2025meta}. Их задача, способ предобучения и критерий качества отличаются от рассматриваемых здесь; эта работа показывает ценность контроля состава предобучающих данных для утверждения о \emph{невиданном} языке. Модель Gamayun демонстрирует отдельное направление исследования небольших многоязычных LLM: двухэтапное предобучение модели с 1,5 млрд параметров с особым вниманием к русскому языку \cite{podolskiy2025gamayun}. В настоящей работе модель с нуля не обучается, поэтому результаты Gamayun не служат ожидаемым значением метрик адаптации.

В сборнике НИТУ «МИСиС» опубликованы работы о классификации русских комментариев с помощью RuBERT \cite{izmaylov2023rubert} и об использовании предобученных моделей в корпоративных чат-ботах \cite{beresnev2023chatbots}. Они дают локальный контекст применения готовых моделей, но не измеряют перенос русской задачи на татарский; их результаты не используются как подтверждение исследуемой гипотезы.

Таким образом, предшествующие работы мотивируют раздельное обучение задачи и языка, контроль токенизации и сравнение компактных обновлений. Они не определяют заранее, какой из вариантов расположения LoRA даст лучший перенос в рассматриваемой постановке.

\section{Постановка задачи}

Пусть $M_0$ — одна открытая базовая авторегрессионная модель с фиксированной ревизией весов и токенизатора. Русская обучающая часть SIB-200 используется для получения общего состояния $M_{src}$, настроенного на тематическую классификацию. Параметры этого этапа выбираются только по русской валидации; минимально приемлемое качество на ней, с предсказанием большинства классов как нижним ориентиром, задаётся до основного эксперимента. Затем каждая языковая адаптация начинает обучение с независимой копии одного и того же $M_{src}$, используя отдельный неразмеченный татарский корпус $C_{tt}$ и функцию потерь предсказания следующего токена. Размеченные татарские примеры при обучении и подборе настроек не используются.

Предмет сравнения — изменение качества на татарском тесте SIB-200 относительно состояния без языковой адаптации. Второй критерий — сохранение качества на русском тесте. До выбора конкретной модели нельзя предполагать, что она слабо знает татарский: предварительный аудит опирается на документацию предобучения, токенизацию и языковую потерю на неразмеченном татарском тексте. Освоение исходной задачи проверяется на русской валидации. Если эти проверки требуют смены модели, кандидат и причина изменения фиксируются до основного сравнения. Татарские метки используются только при итоговой оценке; высокий исходный результат на татарском после её открытия интерпретируется как ограничение выбранного стенда, а не основание для замены модели задним числом.

Основные условия: (1) $M_{src}$ без языковой адаптации; (2) LoRA в нижних блоках; (3) LoRA в верхних блоках; (4) LoRA во всех блоках. Для условий 2 и 3 фиксируются одинаковые набор линейных матриц, число блоков, ранг и масштаб LoRA. Поскольку выбранные блоки одной архитектуры имеют одинаковые размерности, это обеспечивает равное число обучаемых параметров в основной паре. Для условия 4 ранг подбирается к ближайшему достижимому параметрическому бюджету, а расхождение и фактический ранг показываются отдельно. Сопоставление условия 4 с условиями 2–3 интерпретируется как сравнение конфигураций, а не чистый эффект глубины. Обновление только эмбеддингов рассматривается как отдельный контроль с собственным числом параметров. Для моделей со связанными входными эмбеддингами и выходной проекцией эта связь учитывается при подсчёте и интерпретации. Полное дообучение возможно только при достаточном вычислительном бюджете и не входит в обязательное сравнение.

Равенство числа обучаемых параметров не означает равенства вычислительных затрат. Поэтому методы получают одинаковое число обработанных татарских токенов в основной серии, а шаги, время и пиковое потребление памяти фиксируются отдельно. Ранги и объёмы данных могут варьироваться только в дополнительной серии после основного сравнения расположения адаптеров.

\section{Методика эксперимента}

\subsection{Данные и защита от пересечений}

SIB-200 содержит семь тематических классов; в опубликованном разбиении на каждый язык приходится 701 обучающий, 99 валидационных и 204 тестовых предложения \cite[табл.~1]{adelani2023sib}. Поскольку набор получен из параллельных предложений FLORES-200, русские и татарские версии одного исходного предложения имеют общую метку. До обучения проверяется, что исходные предложения и все известные переводы остаются в одном разбиении; в частности, ни одна версия целевого теста не входит в русский train. Для этого используются исходные идентификаторы предложений, а при их отсутствии — таблица соответствия и нормализованные тексты. FLORES-200 не считается независимым корпусом языковой адаптации.

Татарский неразмеченный корпус выбирается с проверенной лицензией и документированным происхождением. Его документы сначала нормализуются и группируются по точным или близким совпадениям, затем группы разделяются на train, validation и test. Из всех частей удаляются обнаруженные совпадения с предложениями SIB-200/FLORES-200; правила нормализации и порог поиска близких копий фиксируются. Train используется для адаптации, validation — только для выбора её настроек, test — для итоговой языковой потери. Межъязыковые совпадения в большом неразмеченном корпусе выявляются не полностью; это остаточный риск утечки, который нужно описать по итогам аудита.

Для выбранного токенизатора до адаптации измеряются число токенов на предложение и распределение длин на русском и татарском текстах. Это позволяет описать эффективность кодирования обоих языков, но само по себе не отделяет её причинный вклад в качество. Ревизии датасетов, правила очистки, хеши исходных файлов и итоговые списки идентификаторов сохраняются вместе с кодом эксперимента.

\subsection{Обучение и выбор настроек}

Первый этап — обучение общей модели $M_{src}$ на русских метках. Каждому из семи классов сопоставляется постоянное обозначение, а формат запроса с русской инструкцией и способ оценки обозначений определяются по русской валидации. Предпочтителен набор обозначений, каждое из которых является одним токеном выбранного токенизатора в данной позиции запроса; это проверяется до обучения. Если такого набора нет, правило оценки полных последовательностей обозначений и нормировки по длине фиксируется заранее. На русской и татарской оценке, а также для всех условий адаптации меняется только классифицируемый текст: инструкция, набор обозначений, их порядок и правило выбора класса остаются теми же. Такая схема не требует генерировать названия тем на татарском, но русская инструкция и обозначения могут влиять на перенос.

На втором этапе независимые копии $M_{src}$ обучаются предсказанию следующего токена на $C_{tt}^{train}$. В основной серии используются одинаковые последовательность данных, число обработанных токенов и правила остановки; размер батча и длина последовательности фиксируются до запуска. Для каждого метода выделяется одинаковое число пробных конфигураций скорости обучения, выбираемых по потере на $C_{tt}^{validation}$ при заранее заданном условии сохранения качества на русской валидации. Точная сетка, оптимизатор и бюджет будут заданы до запуска и не меняются после просмотра татарского теста. Для каждого условия сохраняются конфигурация, seed, исходный checkpoint, версия программной среды, точность вычислений, число шагов, время и пиковая память. Татарские метки открываются только для итоговой оценки; выбор настроек и модели по ним запрещён.

\subsection{Оценка}

Основные метрики — macro-F1 и accuracy на татарском тесте SIB-200. Те же метрики на русском тесте показывают, сохранилась ли исходная задача. Для каждого варианта отдельно показываются абсолютные значения и разности относительно общего $M_{src}$ на тех же тестовых предложениях. До языковой адаптации качество $M_{src}$ на русской валидации сравнивается с заранее заданным порогом и предсказанием большинства классов; если порог не достигнут, основной эксперимент откладывается, а причина смены модели или схемы настройки записывается. Отдельно измеряется средняя потеря предсказания следующего токена на $C_{tt}^{test}$ с одним фиксированным токенизатором. Снижение этой потери само по себе не доказывает улучшения классификации. Отчёт содержит число обучаемых и всех параметров, объём использованных данных в токенах, время и пиковую память.

Если ресурс позволит, каждый вариант выполняется с несколькими seed. При представлении интервалов неопределённости будет указано, какие источники разброса они охватывают: инициализацию адаптера, порядок данных и/или выборку тестовых предложений. При повторной выборке теста русская и татарская версии одного исходного предложения сохраняются парой. Сравнение по одному запуску будет названо пилотным, без утверждений о статистической устойчивости. Основной график покажет изменение татарского и русского качества по условиям адаптации; зависимости от числа параметров и объёма татарских текстов будут строиться только при проведении соответствующих дополнительных серий.

\section{Ограничения постановки}

Публичный состав предобучающих данных модели обычно известен неполно. Даже отсутствие татарского в описании корпуса не гарантирует, что модель не встречала этот язык или предложения FLORES-200. Эксперимент на одной модели, одном языке и одной задаче не устанавливает универсального порядка предпочтения слоёв для всех моделей и задач. Небольшие и несбалансированные классы SIB-200 \cite[табл.~1]{adelani2023sib} ограничивают точность оценки macro-F1; переводное происхождение предложений ограничивает перенос выводов на естественные татарские тексты. При близком параметрическом бюджете варианты могут различаться вычислительной стоимостью, а сравнение LoRA во всех блоках с локальными вариантами дополнительно смешивает глубину и ранг. Выбор фиксированной русской инструкции и обозначений классов тоже может влиять на результат. Кроме того, SIB-200 проверяет тематическую классификацию, а не общее владение татарским языком, рассуждение или наличие языконезависимого «семантического ядра».

До проведения эксперимента работа представляет постановку и методику, а не доказательство эффективности адаптации. После запусков выводы будут ограничены выбранными моделью, корпусом и протоколом, включая возможное отсутствие улучшения.

\begin{thebibliography}{99}
\bibitem{wendler2024llamas} Wendler C., Veselovsky V., Monea G., West R. Do Llamas Work in English? On the Latent Language of Multilingual Transformers. 2024. arXiv:2402.10588. \url{https://arxiv.org/abs/2402.10588}
\bibitem{aghajanyan2020intrinsic} Aghajanyan A., Zettlemoyer L., Gupta S. Intrinsic Dimensionality Explains the Effectiveness of Language Model Fine-Tuning. 2020. arXiv:2012.13255. \url{https://arxiv.org/abs/2012.13255}
\bibitem{hu2021lora} Hu E. J. et al. LoRA: Low-Rank Adaptation of Large Language Models. 2021. arXiv:2106.09685. \url{https://arxiv.org/abs/2106.09685}
\bibitem{hameed2024rosa} Abdel Hameed M. G., Milios A., Reddy S., Rabusseau G. ROSA: Random Subspace Adaptation for Efficient Fine-Tuning. 2024. arXiv:2407.07802. \url{https://arxiv.org/abs/2407.07802}
\bibitem{pfeiffer2020mad} Pfeiffer J., Vulić I., Gurevych I., Ruder S. MAD-X: An Adapter-Based Framework for Multi-Task Cross-Lingual Transfer. 2020. arXiv:2005.00052. \url{https://arxiv.org/abs/2005.00052}
\bibitem{csaki2023efficiently} Csaki Z., Pawakapan P., Thakker U., Xu Q. Efficiently Adapting Pretrained Language Models to New Languages. 2023. arXiv:2311.05741. \url{https://arxiv.org/abs/2311.05741}
\bibitem{adelani2023sib} Adelani D. I. et al. SIB-200: A Simple, Inclusive, and Big Evaluation Dataset for Topic Classification in 200+ Languages and Dialects. EACL, 2024. arXiv:2309.07445. \url{https://arxiv.org/abs/2309.07445}
\bibitem{africa2025meta} Africa D. D. et al. Meta-Pretraining for Zero-Shot Cross-Lingual Named Entity Recognition in Low-Resource Philippine Languages. 2025. arXiv:2509.02160. \url{https://arxiv.org/abs/2509.02160}
\bibitem{podolskiy2025gamayun} Podolskiy A. et al. Gamayun's Path to Multilingual Mastery: Cost-Efficient Training of a 1.5B-Parameter LLM. 2025. arXiv:2512.21580. \url{https://arxiv.org/abs/2512.21580}
\bibitem{izmaylov2023rubert} Измайлов Л. С., Устинов А. М. Эмоциональный окрас комментариев на русском языке с помощью RuBERT моделей // Искусственный интеллект в промышленных, коммерческих, медицинских и финансовых приложениях: сб. ст. науч.-техн. семинара студентов. Вып. 1 / под ред. А. Р. Ефимова. М.: НИТУ «МИСиС», 2023. С. 3–8. \url{https://sadekov.su/Articles/kik_2023.pdf#page=5}
\bibitem{beresnev2023chatbots} Береснев Д. В. Анализ подходов к использованию предобученных моделей в разработке корпоративных чат-ботов // Искусственный интеллект в промышленных, коммерческих, медицинских и финансовых приложениях: сб. ст. науч.-техн. семинара студентов. Вып. 1 / под ред. А. Р. Ефимова. М.: НИТУ «МИСиС», 2023. С. 68–77. \url{https://sadekov.su/Articles/kik_2023.pdf#page=70}
\end{thebibliography}
\end{document}
